\SourceLeaf{264}
\Article{13}
Postquam productum e duobus factoribus evolvere docuimus, evolutio producti e factoribus pluribus nulli difficultati obnoxia erit. Producto $[r,R^\mu]\times[r,R^\nu]$ ad formam $A[1,R^{\mu+\nu}]+B[r,R^{\mu+\nu}]$ reducto, patet, si accedat factor tertius $[r,R^\pi]$, productum fieri $=C[1,R^{\mu+\nu+\pi}]+D[r,R^{\mu+\nu+\pi}]$ statuendo
\[
[r,R^{\mu+\nu}][r,R^\pi]=c[1,R^{\mu+\nu+\pi}]+d[r,R^{\mu+\nu+\pi}]
\]
atque
\[
\begin{aligned}
C&=Bc\\
D&=Bd+A\{1+R^{\mu+\nu}+R^{2\mu+2\nu}+\text{etc.}+R^{(n-2)(\mu+\nu)}\}
\end{aligned}
\]
Hinc potest $[r,R]^\lambda$ facile ad formam $A[1,R^\lambda]+B[r,R^\lambda]$ reduci poterit.

Exempli caussa evolvemus potestates functionis $[r,R]$ pro $n=11$, $\beta=5$, ubi statuemus $g=2$. Hinc respondebunt
\[
\begin{array}{r@{\quad}r@{\pdot}r@{\pdot}r@{\pdot}r@{\pdot}r@{\pdot}r@{\pdot}r@{\pdot}r@{\pdot}r@{\pdot}r}
\text{numeris}&1&2&3&4&5&6&7&8&9&10\\
\text{indices}&0&1&8&2&4&9&7&3&6&5
\end{array}
\]
Habemus itaque ad evolutionem quadrati $[r,R]^2$ secundum formulam I art. 11:
\[
\mu=1,\quad\nu=1
\]
\[
\begin{array}{l@{\quad}r@{\pdot}r@{\pdot}r@{\pdot}r@{\pdot}r@{\pdot}r@{\pdot}r@{\pdot}r@{\pdot}r}
\text{valores ipsius }x\ \ldots&1&2&3&4&5&6&7&8&9\\
\operatorname{ind}x\ \ldots&0&1&8&2&4&9&7&3&6\\
2\operatorname{ind}(x+1)\ \ldots&2&16&4&8&18&14&6&12&10\\
\multicolumn{10}{l}{\text{Res. min. ipsius }\operatorname{ind}x-2\operatorname{ind}(x+1)}\\
\quad\text{secundum modulum }5\ \ldots&3&0&4&4&1&0&1&1&1
\end{array}
\]
unde deducimus
\SourceLeaf{265}
\[
\Omega=[r,R^2]\times\{2+4R+R^3+2R^4\}
\]
atque
\[
1^0.\qquad[r,R]^2=[1,R^2]+[r,R^2]\times\{2+4R+R^3+2R^4\}
\]
Eadem expressio resultat ex formula III art. 11 scilicet
\[
\begin{array}{l@{\quad}r@{\pdot}r@{\pdot}r@{\pdot}r@{\pdot}r@{\pdot}r@{\pdot}r@{\pdot}r@{\pdot}r}
\text{valores ipsius }z\ \ldots&2&3&4&5&6&7&8&9&10\\
\operatorname{ind}z\ \ldots&1&8&2&4&9&7&3&6&5\\
\operatorname{ind}(n+1-z)\ \ldots&5&6&3&7&9&4&2&8&1\\
\multicolumn{10}{l}{\text{resid. min. ipsius }\operatorname{ind}z+\operatorname{ind}(n+1-z)}\\
\quad\text{secundum modulum }5\ \ldots&1&4&0&1&3&1&0&4&1
\end{array}
\]

Prorsus simili modo invenitur
\[
\begin{aligned}
2^0.\qquad[r,R^2]\pdot[r,R]&=[1,R^3]+[r,R^3]\times\{2+R+4R^2+2R^3\}\\
3^0.\qquad[r,R^3]\pdot[r,R]&=[1,R^4]+[r,R^4]\times\{2+4R+R^3+2R^4\}
\end{aligned}
\]
Denique fit
\[
4^0.\qquad[r,R^4]\pdot[r,R]=[1,1]+[r,1]\times\{1+2R+2R^2+2R^3+2R^4\}
\]
Hinc multiplicando aequationem $1^0$ per $[r,R]$ et substituendo pro $[r,R^2]\pdot[r,R]$ valorem suum ex $2^0$, nec non
\[
[1,R^2]\pdot[r,R]=[r,R^3]\pdot\{2+2R+2R^2+2R^3+2R^4\}
\]
deducimus
\[
\begin{aligned}
[r,R]^3&=[1,R^3]\times\{2+4R+R^3+2R^4\}\\
&\quad+[r,R^3]\times\{12+22R+18R^2+24R^3+15R^4\}
\end{aligned}
\]
et simili modo
\[
\begin{aligned}
[r,R]^4&=[1,R^4]\times\{12+22R+18R^2+24R^3+15R^4\}\\
&\quad+[r,R^4]\times\{164+170R+205R^2+180R^3+190R^4\}\\
\relax[r,R]^5&=[1,1]\times\{164+170R+205R^2+180R^3+190R^4\}\\
&\quad+[r,1]\times\{1836+1830R+1795R^2+1820R^3+1810R^4\}\\
&=1640+1700R+2050R^2+1800R^3+1900R^4\\
&\quad+(1836+1830R+1795R^2+1820R^3+1810R^4)(s-1)\\
&=918Ss-98S-(6R+41R^2+16R^3+26R^4)s\\
&\quad+66R+451R^2+176R^3+286R^4
\end{aligned}
\]

\SourceLeaf{266}
\Article{14}
Calculus in praecc.\textsuperscript{'} ita absolutus, ut ad omnes valores ipsius $r$ ipsiusque $R$ extendi possit, notabiliter contrahitur, si ipsam $R$ statim ab initio tamquam radicem propriam aequationis $x^\beta-1=0$ consideramus. Hacce suppositione productum $[r,R^\mu]\times[r,R^\nu]$ reducetur ad formam $B[r,R^{\mu+\nu}]$, quoties $\mu+\nu$ per $\beta$ non est divisibilis; quando vero $\mu+\nu$ per $\beta$ divisibilis est, illud productum fit $=(n-1)R^{\frac12(n-1)\mu}+[r,1]\Sigma R^{\mu\operatorname{ind}x}$, substituendo pro $\operatorname{ind}x$ omnes numeros $0,1,2,3\dotsfour n-2$ excepto hoc $\frac12(n-1)$. Hinc facile colligitur (si $\mu$ et proin etiam $\nu$ per $\beta$ non est divisibilis), in hoc casu esse
\[
[r,R^\mu]\pdot[r,R^\nu]=R^{\frac12(n-1)\mu}\{n-1-[r,1]\}
\]
adeoque $=0$ pro $r=1$, et $=nR^{\frac12(n-1)\mu}$ pro quovis alio valore ipsius $r$. Ceterum quum $R^{\frac12(n-1)\mu}$ fiat $=+1$, vel $=-1$, prout $\frac{n-1}\beta\pdot\mu$ est numerus par vel impar, productum nostrum fit in casu priori $=n$, in posteriori $=-n$.

Hinc porro sequitur, statui posse
\[
\begin{aligned}
[r,R]^2&=A'[r,R^2]\\
\relax[r,R^2]\pdot[r,R]&=A''[r,R^3]\\
\relax[r,R^3]\pdot[r,R]&=A'''[r,R^4]
\end{aligned}
\]
etc. usque ad
\[
[r,R^{\beta-2}]\pdot[r,R]=A^{(\beta-2)}[r,R^{\beta-1}]
\]
unde habemus
\[
\begin{aligned}
[r,R]^2&=A'[r,R^2]\\
\relax[r,R]^3&=A'A''[r,R^3]\\
\relax[r,R]^4&=A'A''A'''[r,R^4]
\end{aligned}
\]
etc. Denique
\[
[r,R]^\beta=\pm nA'A''A'''\dots A^{(\beta-2)}
\]
ubi signum superius vel inferius accipiendum est, prout $\frac{n-1}\beta$ par est vel impar.

Patet itaque, postquam valor ipsius $[r,R]$ inventus fuerit, functiones reliquas
\[
[r,R^2]=\frac{[r,R]^2}{A'},\quad[r,R^3]=\frac{[r,R]^3}{A'A''}\quad\text{etc.}
\]
hic multo expeditius determinari posse, quam in casibus iis, ubi $\alpha$ non est $=1$,
\SourceLeaf{267}
ut iam in \emph{Disq. Ar.} (Art. 360, \textsc{iii}) monuimus. Per considerationem uberiorem indolis functionum $A',A''$ etc. hae operationes adhuc magis facilitantur.

\Article{15}
In art. 9 ostendimus, valorem functionis $[r,R]$ reduci posse ad formam $\surd n(\cos f+i\sin f)$, eodemque modo functiones $[r,R^2]$, $[r,R^3]$ etc. usque ad $[r,R^{\beta-1}]$ ad similem formam reduci poterunt. Statuamus
\[
\begin{aligned}
[r,R]&=\surd n(\cos f'+i\sin f')\\
\relax[r,R^2]&=\surd n(\cos f''+i\sin f'')\\
\relax[r,R^3]&=\surd n(\cos f'''+i\sin f''')\\
&\text{etc.}
\end{aligned}
\]
eritque
\[
\begin{aligned}
A'&=\surd n(\cos(2f'-f'')+i\sin(2f'-f''))\\
A''&=\surd n(\cos(f'+f''-f''')+i\sin(f'+f''-f'''))\\
A'''&=\surd n(\cos(f'+f'''-f'''')+i\sin(f'+f'''-f''''))\\
&\text{etc.}
\end{aligned}
\]
Hinc patet, si functiones $A',A'',A'''$ etc. reducantur ad formas
\[
\begin{aligned}
A'&=a'(\cos b'+i\sin b')\\
A''&=a''(\cos b''+i\sin b'')\\
A'''&=a'''(\cos b'''+i\sin b''')\\
&\text{etc.}
\end{aligned}
\]
et quidem ita, ut omnes $a',a'',a'''$ etc. sint positivi, fore
\[
\begin{aligned}
a'&=a''=a'''\text{ etc.}=\surd n\\
f'&=\frac1\beta(b'+b''+b'''+\text{etc.}+b^{(\beta-2)})
\end{aligned}
\]
si fuerit $\frac{n-1}\beta$ par, vel
\[
f'=\frac1\beta(180^0+b'+b''+\text{etc.}+b^{(\beta-2)})
\]
si fuerit $\frac{n-1}\beta$ impar, ac dein
\[
\begin{aligned}
[r,R]&=\surd n(\cos f'+i\sin f')\\
\relax[r,R^2]&=\surd n(\cos(2f'-b')+i\sin(2f'-b'))\\
\relax[r,R^3]&=\surd n(\cos(3f'-b'-b'')+i\sin(3f'-b'-b''))\\
&\text{etc.}
\end{aligned}
\]
\SourceLeaf{268}
denique erit per formulas art. 8
\[
\begin{aligned}
\left(\frac{n-1}\beta,1\right)&=-\frac1\beta+\frac{\surd n}\beta\{\cos f''+\cos(2f'-b')+\cos(3f'-b'-b'')+\text{etc.}\\
&\qquad+\cos((\beta-1)f'-b'-b''-b'''-\text{etc.}-b^{(\beta-2)})\}\\
&\quad+\frac{i\surd n}\beta\{\sin f'+\sin(2f'-b')+\sin(3f'-b'-b'')+\text{etc.}\}
\end{aligned}
\]
et perinde prodeunt valores functionum $\left(\frac{n-1}\beta,g\right)$, $\left(\frac{n-1}\beta,g^2\right)$, $\left(\frac{n-1}\beta,g^3\right)$ etc., si in hac formula pro $f'$ resp. substituitur $f'-\frac{360^0k}\beta$, $f'-2\frac{360^0k}\beta$, $f'-3\frac{360^0k}\beta$ etc., supponendo $R=\cos\frac{360^0k}\beta+i\sin\frac{360^0k}\beta$.

\Article{16}
Simplificatio nova ex observatione sequente petitur. Quum per art. 14 fiat
\[
\pm[r,R][r,R^{\beta-1}]=[r,R^2][r,R^{\beta-2}]=\pm[r,R^3][r,R^{\beta-3}]\text{ etc.}=n
\]
accipiendo [in producto primo, tertio etc.] signum superius vel inferius, prout $\frac{n-1}\beta$ par est vel impar, esse debebit in casu priori
\[
\cos(f'+f^{(\beta-1)})=\cos(f''+f^{(\beta-2)})=\cos(f'''+f^{(\beta-3)})\text{ etc.}=1
\]
in posteriori
\[
-\cos(f'+f^{(\beta-1)})=\cos(f''+f^{(\beta-2)})=-\cos(f'''+f^{(\beta-3)})\text{ etc.}=1
\]
et in utroque casu
\[
\sin(f'+f^{(\beta-1)})=\sin(f''+f^{(\beta-2)})=\sin(f'''+f^{(\beta-3)})\text{ etc.}=0
\]
Hinc statuere licebit in casu priori
\[
f^{(\beta-1)}=-f',\quad f^{(\beta-2)}=-f'',\quad f^{(\beta-3)}=-f'''\quad\text{etc.}
\]
in posteriori
\[
f^{(\beta-1)}=180^0-f',\quad f^{(\beta-2)}=-f'',\quad f^{(\beta-3)}=180^0-f'''\quad\text{etc.}
\]
hinc vero sequitur, in priori casu esse
\[
\begin{aligned}
b^{(\beta-2)}&=b',&b^{(\beta-3)}&=b'',&b^{(\beta-4)}&=b'''\quad\text{etc.}\\
A^{(\beta-2)}&=A',&A^{(\beta-3)}&=A'',&A^{(\beta-4)}&=A'''\quad\text{etc.}
\end{aligned}
\]
in posteriori vero
\SourceLeaf{269}
\[
\begin{aligned}
b^{(\beta-2)}&=b'-180^0,&b^{(\beta-3)}&=b''+180^0,&b^{(\beta-4)}&=b'''-180^0\quad\text{etc.}\\
A^{(\beta-2)}&=-A',&A^{(\beta-3)}&=-A'',&A^{(\beta-4)}&=-A'''\quad\text{etc.}
\end{aligned}
\]
ita ut multitudo functionum $A',A'',A'''$ etc. ad semissem reducatur. Hinc porro colligitur, in priori casu fore
\[
f'=\frac1\beta(2b'+2b''+\text{etc.}+2b^{(\frac12\beta-1)})
\]
\[
\begin{aligned}
\left(\frac{n-1}\beta,1\right)&=-\frac1\beta+\frac{\surd n}\beta\{2\cos f'+2\cos(2f'-b')\\
&\qquad+2\cos(3f'-b'-b'')+\text{etc.}\\
&\qquad+2\cos((\tfrac12\beta-1)f'-b'-b''-\text{etc.}-b^{(\frac18\beta-2)})\\
&\qquad+\cos(\tfrac12\beta f'-b'-b''-\text{etc.}-b^{(\frac12\beta-1)})\}
\end{aligned}
\]
(ubi terminus ultimus manifesto est $=\cos0=1$) vel
\[
f'=\frac1\beta(2b'+2b''+\text{etc.}+2b^{(\frac12(\beta-3))}+b^{(\frac12(\beta-1))})
\]
\[
\begin{aligned}
\left(\frac{n-1}\beta,1\right)&=-\frac1\beta+\frac{\surd n}\beta\{2\cos f'+2\cos(2f'-b')\\
&\qquad+2\cos(3f'-b'-b'')+\text{etc.}\\
&\qquad+2\cos(\tfrac12(\beta-1)f'-b'-b''-\text{etc.}-b^{(\frac12(\beta-3))})\}
\end{aligned}
\]
prout $\beta$ par est vel impar; et in casu posteriori

\[
f'=\frac1\beta(2b'+2b''+\text{etc.}+2b^{(\frac12\beta-1)})
\]
\[
\begin{aligned}
\left(\frac{n-1}\beta,1\right)&=-\frac1\beta+\frac{\surd n}\beta\{2\cos(2f'-b')+2\cos(4f'-b'-b''-b''')+\text{etc.}\\
&\qquad+2\cos((\tfrac12\beta-2)f'-b'-b''-\text{etc.}-b^{(\frac12\beta-3)})\\
&\qquad+\cos(\tfrac12\beta f'-b'-b''-\text{etc.}-b^{(\frac12\beta-1)})\}\\
&\quad+i\frac{\surd n}\beta\{2\sin f'+2\sin(3f'-b'-b'')+\text{etc.}\\
&\qquad+2\sin((\tfrac12\beta-1)f'-b'-b''-\text{etc.}-b^{(\frac12\beta-2)})\}
\end{aligned}
\]
vel
\[
f'=\frac1\beta(2b'+2b''+\text{etc.}+2b^{(\frac12\beta-1)}+180^0)
\]
\[
\begin{aligned}
\left(\frac{n-1}\beta,1\right)&=-\frac1\beta+\frac{\surd n}\beta\{2\cos(2f'-b')+2\cos(4f'-b'-b''-b''')+\text{etc.}\\
&\qquad+2\cos((\tfrac12\beta-1)f'-b'-b''-\text{etc.}-b^{(\frac12\beta-2)})\}\\
&\quad+i\frac{\surd n}\beta\{2\sin f'+2\sin(3f'-b'-b'')+\text{etc.}\\
&\qquad+2\sin((\tfrac12\beta-2)f'-b'-b''-\text{etc.}-b^{(\frac12\beta-3)})\\
&\qquad+\sin(\tfrac12\beta f'-b'-b''-\text{etc.}-b^{(\frac12\beta-1)})\}
\end{aligned}
\]
prout $\frac12\beta$ par est vel impar. De periodis reliquis $\frac{n-1}\beta$ terminorum eadem valent, quae supra annotavimus. Generaliter itaque hinc concluditur, ad determinationem harum periodorum requiri sectionem circuli integri in $\beta$ partes, a qua
\SourceLeaf{270}
constructio angulorum $b',b'',b'''$ etc. rationaliter pendet, dein divisionem anguli $b'+b''+b'''+\text{etc.}$ in $\beta$ partes, denique radicem quadratam $\surd n$. Quodsi statuitur statim $\beta=\frac12(n-1)$, periodi illae manifesto coincidunt cum duplicatis cosinibus angulorum $\frac{360^0}n$, $2\frac{360^0}n$, $3\frac{360^0}n$ etc. usque ad $\frac12(n-1)\frac{360^0}n$, ita ut divisio circuli in $n$ partes pendeat a divisione circuli integri in $\frac12(n-1)$ partes, divisione anguli, qui illa sectione perfecta construi potest, in $\frac12(n-1)$ partes, atque quantitate radicali $\surd n$. Si usque ad sinus angulorum $\frac{360^0}n$ etc. progredi constitutum est, una operatione amplius opus erit.
